\subsection{Interface Contracts and Workstream Division}
\label{subsec:interface-contracts}

A system of this size is only tractable to build if its seams are chosen
deliberately rather than allowed to congeal around whoever happened to
write the first line of code in each area. We therefore fixed the
interfaces between major components in Week~1, before any production
implementation existed on either side of them, and treated those
interfaces as frozen artefacts that could be changed only by mutual
agreement. The pay-off is that the our two human workstreams, a backend
workstream covering the orchestrator, database, portal, and event
pipeline, and an AI/LLM workstream covering context extraction, vector
search, re-ranking, response generation, and engagement judging 
were able to develop in parallel against typed mocks of one another and
integrate in a single short sitting rather than in a protracted
debugging phase.

The contract is expressed as a small set of Python dataclasses and
asynchronous function signatures. Four data types carry almost the
entire load. A \texttt{Turn} is an immutable record of one message in a
conversation, tagged with its role, timestamp, and the identifier of
any ad that was shown alongside it. A \texttt{ContextObject} is the
structured projection of a user's most recent turn, carrying an
enumerated \texttt{intent}, a list of topics and entities, a sentiment
label, a purchase-signal scalar in $[0,1]$, and an \texttt{ad\_receptivity}
enum whose \texttt{none} value short-circuits the matching pipeline
entirely for sensitive conversations. An \texttt{Ad} is the unit of
inventory, carrying targeting metadata, creative text, brand-safety
tags, a bid in CPM, a remaining budget, and two optional fields that are populated in
sequence by the vector search and the re-ranker. A
\texttt{ResponseObject} is what the orchestrator returns to the
frontend: a response text, a boolean indicating whether an ad was
included, and the identifier of the ad used if so. The enums
(\texttt{Intent}, \texttt{AdReceptivity}) are part of the contract
rather than free-form strings so that a typo on one side of the seam
becomes a type error at import time rather than a silent mismatch at
runtime.

Four asynchronous function signatures then pin down the call graph
between the workstreams. \texttt{extract\_context(message, history)
\(\to\) ContextObject} is owned by the AI workstream and called by the
orchestrator once per user turn. \texttt{match\_ads(context, session\_id)
\(\to\) list[Ad]} is also owned by the AI workstream and encapsulates
the embedding, vector search, and re-ranking stages behind a single
entry point, so that the backend need not know whether matching is
performed by a single model call or a multi-stage pipeline.
\texttt{generate\_response(message, history, ad)\(\to\) ResponseObject}
is the main model call that weaves the selected ad (if any) into a
helpful answer. Finally,
\texttt{judge\_engagement(ad, prev\_assistant\_turn, user\_turn) \(\to\)
EngagementVerdict} is the post-hoc judge described in
Section~\ref{subsubsec:llm-judge-engagement}. Everything that is not in
this list --- database access, Redis session state, event writes,
portal CRUD, budget pacing --- is owned unambiguously by the backend
workstream and is invisible to the AI workstream except through these
four functions.

The division of labour has two consequences worth naming. First, either
side can be replaced wholesale without touching the other: swapping the
matching pipeline for a learned end-to-end ranker, or swapping the
Postgres event store for a streaming pipeline, requires changes only
inside one workstream's boundary. Second, testing becomes local. The
backend workstream tests against a deterministic stub of the AI
functions that returns canned \texttt{ContextObject} and \texttt{Ad}
values; the AI workstream tests against an in-memory fake of the
session store and event writer. Neither side needs the other to be
running in order to validate its own invariants.

\subsection{Quality Assurance and Evaluation Framework}
\label{subsec:qa-framework}

Evaluating a conversational ad system is harder than evaluating either
a chat model or a classical ad ranker in isolation, because the quality
of a served ad is a joint function of its relevance to the turn, its
compatibility with the surrounding assistant prose, and its effect on
the user's subsequent behaviour. We therefore adopt a layered
evaluation framework in which each component is held to a local
correctness criterion and the system as a whole is judged against a
small number of end-to-end metrics that reflect the product's stated
value proposition.

At the component level, three properties are checked continuously.
Context extraction is evaluated for \emph{schema validity} and
\emph{label stability}: the context extractor must return a
\texttt{ContextObject} that parses against the dataclass schema on
every call, and its labels on a fixed regression set of fifty hand-
annotated turns must remain stable across prompt revisions. A failure
of the first property is caught at parse time and triggers a fallback
to a conservative default; a failure of the second is caught by a
nightly job that diffs the extractor's output against the regression
set and opens an alert if agreement drops below a threshold. Vector
search is evaluated for \emph{recall at 20} on a small synthetic set
of query--ad pairs constructed from the inventory's own product
descriptions, which gives a cheap, self-generated benchmark that
automatically tracks changes in the embedding model or the inventory.
The re-ranker is evaluated by \emph{rank correlation} between its
output ordering and the ordering induced by the unmodulated similarity
score, with the expectation that the two should diverge in
interpretable ways --- budget-exhausted ads should fall, high
purchase-signal sessions should see higher-bid ads rise --- and that
the divergence itself is what the unit tests assert.

At the end-to-end level, the system is judged on three metrics.
\emph{Engagement rate}, the fraction of served ads whose subsequent
user turn is flagged as engaged by the judge of
Section~\ref{subsubsec:llm-judge-engagement}, is the headline quality
metric and is the one that the feedback loop optimises against.
\emph{Click-through rate} is reported alongside it as a secondary
metric, mainly for comparability with conventional ad analytics.
\emph{Injection appropriateness} is the rate at which ads are served
into turns whose \texttt{ad\_receptivity} is \texttt{high} or
\texttt{medium} and suppressed on turns whose receptivity is
\texttt{low} or \texttt{none}; this metric exists to catch regressions
in the sensitivity classifier independently of whether the served ads
are themselves good, because an inappropriately-placed ad is a product
failure even if the user happens to engage with it.

We deliberately do not use offline A/B testing as a primary evaluation
tool. The feedback loop described in
Section~\ref{subsubsec:adaptive-optimisation} means that serving
decisions on one turn shape context on subsequent turns within the
same session, which violates the independence assumption that
conventional A/B tests rely on. Instead, we rely on the regression
sets and rank-correlation tests above to catch local regressions, and
on the headline engagement rate to track global drift. A principled
bandit-style evaluation that respects the within-session dependency
structure is sketched in Section~\ref{subsec:future-work} as the
natural successor to the current framework.

\subsection{Error Handling and Graceful Degradation}
\label{subsec:error-handling}

The serving path touches three external systems that can each fail
independently. A Postgres instance with a pgvector extension, a
Redis instance holding session state, and one or more remote LLM
endpoints and any production deployment must answer the question
of what the user sees when one of them is unreachable or slow. Our
guiding principle is that \emph{the chat response is the product},
and that every component in the ad pipeline is expendable in the
service of delivering a timely, coherent answer to the user. No
failure in the ad pipeline should ever degrade the conversational
experience, and no ad should ever be served on the back of partial or
inconsistent data.

This principle is realised as a small number of explicit fallback
paths, each triggered by a specific class of failure. A failure of the
context extractor causes the orchestrator to substitute a
default \texttt{ContextObject} whose \texttt{ad\_receptivity} is set
to \texttt{none}. This value is the one the matching pipeline
interprets as ``do not serve an ad on this turn,'' so the effect of
the failure is silent suppression: the user receives a normal
assistant response with no sponsored content, and an error event is
written to the log for later inspection. A failure of the vector
search, whether due to a Postgres outage or an empty candidate set,
is treated identically: the matching function returns an empty list,
the orchestrator observes that no candidates are available, and the
response-generation call is issued without any ad in its prompt. A
failure of Redis is handled by treating the session as if it were
new; the current turn is extracted and matched in isolation, and the
adaptive within-session loop is skipped for that turn only. In all
three cases the system degrades to a strictly weaker but still
correct version of itself, rather than surfacing an error to the
user or blocking the response.

The main-model response generation call is the one external
dependency that cannot be silently elided, because there is no
response without it. For this call we adopt a bounded-retry policy
with exponential backoff and a hard deadline of a few seconds,
beyond which the orchestrator returns a generic error message to the
frontend and logs the incident. Ad-related state is never committed
for a turn whose response generation failed: no impression event is
written, no session state is updated, and no budget is debited.
This preserves the invariant that every impression in the events
table corresponds to an ad the user actually saw.

A second class of failures, subtler than outright outages, concerns
partial inconsistency between components. The portal, for instance,
can mark a campaign as paused between the moment an ad was selected
by the re-ranker and the moment the response is generated; a
campaign's daily budget can be exhausted by a concurrent request in
the same window; an ad's embedding can be regenerated in place after
the candidate list was assembled. We address these races with a
final-check pattern: immediately before the response-generation
call, the orchestrator re-reads the selected ad's status, budget,
and version from the database inside a single short transaction and
aborts the injection if any invariant has been violated since
selection. The cost of this re-read is one indexed point lookup per
turn, which is negligible next to the model call it precedes, and
the benefit is that the events table never contains an impression
for an ad that was paused, out of budget, or stale at the instant
of serving. Taken together, these mechanisms give the system a
clear failure contract: ads may silently disappear, but the
conversational experience does not, and the analytics layer never
reports on events that did not faithfully happen.